% Bibliography entries for task 2; insert in the master thebibliography.
\bibitem{NXLBorger} J. Borger, \emph{$\Lambda$-rings and the field with one element},
arXiv:0906.3146. Supplied source
\path{LIT_Borger_Lambda_rings_and_the_field_with_one_element/lrfoe13.tex};
primary PDF edition dated 26 November 2024, pp.~2 and~6 (definition and flat
Frobenius-lift criterion), \S6 opening and Theorem~6.1, pp.~19--20
(abelian finite-coefficient Galois action), Corollary~6.2 and equation~(6.2.1),
p.~20 (potential integral-Tate splitting). Supplied source labels
\texttt{thm:lambda-varieties-are-abelian} and
\texttt{cor:lambda-varieties-are-Artin-Tate}.
\url{https://arxiv.org/pdf/0906.3146}.

\bibitem{NXLScholze} P. Scholze, \emph{Perfectoid spaces},
Publ. Math. Inst. Hautes \'Etudes Sci. \textbf{116} (2012), 245--313.
The primary author PDF used here is the 19 November 2011 version:
Definition~5.1(ii), Theorem~5.2, and Lemma~5.3, pp.~21--23;
Proposition~5.13(i) and Remark~5.14, p.~25;
Lemma~6.2 and Theorem~6.3, p.~30;
Theorem~7.12, p.~43;
Definition~7.14, pp.~43--44, and Theorem~7.17, Corollaries~7.18--7.19,
p.~44;
Theorem~9.7 (citing Huber, Theorem~3.8.1), p.~49.
\url{https://people.mpim-bonn.mpg.de/scholze/PerfectoidSpaces.pdf}.

\bibitem{NXLHuber} R. Huber, \emph{A generalization of formal schemes and rigid
analytic varieties}, Math. Z. \textbf{217} (1994), 513--551,
Remark~4.6(iv)(d), printed p.~544. The hypotheses are locally noetherian
schemes and a proper morphism; the characteristic-$p$ application in NX-L02
is descended to $\mathbb F_p[[u]]$ before base change.
\url{https://virtualmath1.stanford.edu/~conrad/Perfseminar/refs/Huberformalrigid.pdf}.

\bibitem{NXLBhatt} B. Bhatt, \emph{Lecture II: $\delta$-rings},
Columbia prismatic-cohomology lectures, Proposition~3.2, pp.~5--6
(unique strict-$p$ lift of a perfect $\mathbb F_p$-algebra).
The primary PDF text was retrieved; the direct local download returned 403.
The proof can use Scholze, Remark~5.14, alone.
\url{https://websites.umich.edu/~bhattb/teaching/prismatic-columbia/lecture2-delta-rings.pdf}.

\bibitem{NXLSilverman} J. H. Silverman, \emph{The Arithmetic of Elliptic Curves},
second edition, Graduate Texts in Mathematics 106, Springer, 2009,
Theorem~V.2.4, printed p.~143 (elliptic-curve zeta numerator).
The explicit trace $a_5=-2$ in NX-L01 is recomputed directly.
\url{https://www.math.ens.psl.eu/~obenoist/refs/Silverman.pdf}.

\bibitem{NXLSGA4} M. Artin, A. Grothendieck, J.-L. Verdier, eds.,
\emph{Th\'eorie des topos et cohomologie \'etale des sch\'emas (SGA~4)},
Expos\'e XII, Theorem~5.1 and Corollary~5.2 (proper base change);
Expos\'e XVI, Theorem~2.1 (smooth base change);
Expos\'e XVII, \S\S5.1--5.3 (compact supports).
Exact locators carried from the supplied predecessor; no new scan
verification or re-proof of these source theorems is claimed.

\bibitem{NXLSGA5} A. Grothendieck et al.,
\emph{Cohomologie $\ell$-adique et fonctions $L$ (SGA~5)},
Expos\'e XV, \S3, no.~2, unnumbered theorem, equation~(1)
(finite-field trace formula). Locator carried from ZETA\_TWO~ZT04.

\bibitem{NXLSGA6} A. Grothendieck et al.,
\emph{Th\'eorie des intersections et th\'eor\`eme de Riemann--Roch (SGA~6)},
Expos\'e V, \S2 and \S7, equation~(7.1.1)
(exterior powers, symmetric powers and determinant formalism).
Locator carried from ZETA\_TWO~ZT08 and SV01.

\bibitem{NXLPredecessors} \emph{ZETA\_TWO}, final supplied run
\texttt{RETURN\_ZETA\_TWO\_final\_20261121T0032Z}:
ZT01 (abelian Artin--Tate transfer under A), ZT04 (finite-factor correction),
ZT05 (omitted-prime zeros), ZT08 and SV01 (determinant/$\lambda$ formalism);
\texttt{INTAKE\_ZETA\_TWO\_20261121.md}, \S5.2.
\emph{DP\_MULT}, supplied checkpoint
\texttt{RETURN\_DP\_MULT\_checkpoint\_03\_20261121T0205Z}, DP11:
the binomial/identity-Adams special case, not a criterion for every
$\Lambda$-ring to be binomial.

\bibitem{NXLArchive} Grothendieck archive locators carried from the supplied
ZETA\_TWO \emph{IMPLICATIONS}, \S5 placement table: cote~60,
scans~3 and~5 (exterior and symmetric powers), scan~7, archival p.~6
(reciprocal characteristic polynomial), scan~9, archival p.~8
(Weil and dual/Tate normalization). DP\_MULT carries cote~153,
PDF p.~3/archive p.~2 (binomial polynomials). These locators were not
newly transcribed or visually inspected in this run; none is represented
as a proof of NX-L02 or of bad-prime splitting.

\bibitem{NXLAlpoge} L. Alp\"oge, \emph{Remark}, supplied four-page note,
20 November 2026, step~20, p.~4; supplied
\texttt{ALPOGE\_Main.lean}, declarations
\texttt{LFunction\_eq\_zero\_re} and
\texttt{LFunction\_ne\_zero\_of\_half\_lt\_re}, with definitions in
\texttt{ALPOGE\_Statements.lean}. Used as stipulated theorem A only
in NX-L06; NX-L01--NX-L05 and NX-L07 use neither A nor Masullo's M.
\url{https://alpo.ge/remark.pdf}.
